% ============================================================================
%  Chapter 9 -- Elements of Boolean Algebra, Part 2              (Lecture 09)
% ============================================================================
\chapter{Boolean Algebra: Functions, Canonical Forms and Minimisation}
\label{ch:boole2}

\epigraphline{Oh dear! Oh dear! I shall be too late!}
  {Lewis Carroll, \emph{Alice's Adventures in Wonderland}, 1865}

\begin{chapgoals}
By the end of this chapter you will be able to
\begin{itemize}
\item define literal, clause, factor, minterm and maxterm;
\item say how many Boolean functions of $n$ variables exist, and name those of two variables;
\item write a function in canonical sum-of-products and product-of-sums form from its truth table;
\item minimise a function of up to four variables with a Karnaugh map;
\item derive the half adder and the full adder, and explain how thirty-two of them execute \ins{add}.
\end{itemize}
\end{chapgoals}
% Generated by tools/textbook/export_listings.js from tools/decks/mastery.js and tools/concepts/registry.json: do not edit.
\begin{chapstatus}{Learning status of this chapter}
\masterychip{now}{UNDERSTAND NOW}\enspace the adder; Boolean functions; canonical forms; minimisation, Karnaugh maps.\par\smallskip
\masterychip{use}{USE NOW}\enspace the complete program, the PRG32 tools (used, not studied, in this course).\par\smallskip
\masterychip{later}{UNDERSTAND LATER}\enspace nothing.
\end{chapstatus}
\begin{chapstatus}{Roadmap of this chapter}
By the end of the chapter you can:
\begin{enumerate}[nosep,leftmargin=2em]
\item write a function in canonical form from its truth table
\item minimise a function of up to four variables with a Karnaugh map
\item build a full adder and chain it into an n-bit adder
\end{enumerate}
It builds on: only the setup of Chapter~\ref{ch:setup}. Next: Chapter~\ref{ch:electronics}.
\end{chapstatus}


% ----------------------------------------------------------------------------
\section{Warm-up on the machine}

Two registers hold one bit each, 0 or 1. Their sum can be 0, 1 or 2, and
therefore needs two bits. Listing~\ref{lst:b2:warm} computes it without any
arithmetic instruction.

\asmfile[caption={One-bit addition without \ins{add}.},label=lst:b2:warm]{code/l09/c03.S}

Two logic instructions add two bits. The comment of the listing is a truth
table with two outputs, and the pair of gates that it describes is called a
\term{half adder}. Chapter~\ref{ch:integers} claimed that an adder is all a
processor needs for arithmetic: this is its first bit. But how are the gates
found when only the table is given? That is the method of this chapter:
truth table, canonical form, minimisation.

% ----------------------------------------------------------------------------
\section{Boolean functions}

\subsection{The vocabulary of expressions}

\begin{definition}{Literal, clause, factor}
A \term{literal} is a Boolean variable, or its complement, as it appears in an
expression: $a$, $\nott{b}$, $x$. A \term{clause}, or \term{product term}, is
the logical product of two or more literals: $a\,\nott{b}\,x$. A
\term{factor}, or \term{sum term}, is the logical sum of two or more literals:
$(a+\nott{x})$, $(b+c)$.
\end{definition}

Clauses and factors with more than one literal are functions of level~1: one
layer of gates. Each occurrence of a variable counts as a literal, and the
number of literals measures the cost of an expression.

\subsection{How many functions of $n$ variables?}

\begin{definition}{Boolean function}
A \term{Boolean function} of $n$ variables is a map $f:B^{n}\to B$, with
$B=\{0,1\}$: it associates one binary output with every combination of $n$
binary inputs.
\end{definition}

\begin{theorem}{Number of Boolean functions}
There are exactly $2^{2^{n}}$ distinct Boolean functions of $n$ variables.
\end{theorem}

\noindent\emph{Proof.} There are $2^{n}$ combinations of the inputs, and a
function is determined by the choice of an output, 0 or 1, for each of them:
it is a string of $2^{n}$ bits, and there are $2^{2^{n}}$ such strings.
\hfill$\square$

\smallskip
The number is finite, and it grows quickly: 4 functions of one variable, 16 of
two, 256 of three, 65\,536 of four, and more than four thousand million of
five. Boolean functions are isomorphic to digital circuits: every circuit can
be written as an expression, and every expression can be built as a circuit.

\subsection{The sixteen functions of two variables}

Table~\ref{tab:b2:sixteen} lists all the functions of two variables. The
index of each function is its column of outputs read as a binary number. With
one variable the four functions are the constant 0, $x$, $\nott{x}$ and the
constant 1.

\begin{table}[htb]
\centering\small
\caption{The sixteen functions of two variables. The columns 00, 01, 10, 11
give the output for each pair of inputs $x\,y$.}\label{tab:b2:sixteen}
\begin{tabular}{@{}lccccl@{\qquad}lccccl@{}}
\toprule
$f$ & 00 & 01 & 10 & 11 & \textbf{name} & $f$ & 00 & 01 & 10 & 11 & \textbf{name}\\
\midrule
$f_0$ & 0&0&0&0 & constant 0 & $f_8$ & 1&0&0&0 & \textbf{NOR}\\
$f_1$ & 0&0&0&1 & \textbf{AND}: $x\,y$ & $f_9$ & 1&0&0&1 & \textbf{XNOR} (equivalence)\\
$f_2$ & 0&0&1&0 & $x\,\nott{y}$ (inhibition) & $f_{10}$ & 1&0&1&0 & NOT $y$\\
$f_3$ & 0&0&1&1 & $x$ (identity) & $f_{11}$ & 1&0&1&1 & $x+\nott{y}$ (implication $y\to x$)\\
$f_4$ & 0&1&0&0 & $\nott{x}\,y$ (inhibition) & $f_{12}$ & 1&1&0&0 & NOT $x$\\
$f_5$ & 0&1&0&1 & $y$ (identity) & $f_{13}$ & 1&1&0&1 & $\nott{x}+y$ (implication $x\to y$)\\
$f_6$ & 0&1&1&0 & \textbf{XOR} & $f_{14}$ & 1&1&1&0 & \textbf{NAND}\\
$f_7$ & 0&1&1&1 & \textbf{OR}: $x+y$ & $f_{15}$ & 1&1&1&1 & constant 1\\
\bottomrule
\end{tabular}
\end{table}

\subsection{Functionally complete sets}

Any Boolean function can be realised using only AND, OR and NOT; for example
$x\oplus y=\nott{x}\,y+x\,\nott{y}$. (Section~\ref{sec:b2:canonical} proves
the general statement by construction.) The set $\{\text{AND},\text{OR},
\text{NOT}\}$ is therefore called \term{functionally complete}.

By De Morgan, AND can be built from OR and NOT, and OR from AND and NOT: the
sets $\{\text{AND},\text{NOT}\}$ and $\{\text{OR},\text{NOT}\}$ are complete
too. And a single gate is enough.

\begin{theorem}{Universal gates}
NAND alone is functionally complete, and so is NOR alone.
\end{theorem}

\noindent\emph{Proof} for NAND, writing $x\uparrow y=\nott{xy}$:
\[
\nott{x}=x\uparrow x,\qquad
x\,y=(x\uparrow y)\uparrow(x\uparrow y),\qquad
x+y=(x\uparrow x)\uparrow(y\uparrow y),
\]
the last by De Morgan. The proof for NOR is dual. \hfill$\square$

\smallskip
A single type of gate is enough to build any digital circuit, a whole
processor included. Real chips are largely built from NAND and NOR because
these are the gates that need the fewest transistors
(Chapter~\ref{ch:electronics}).

\subsection{Truth tables}

A Boolean function can be defined by its truth table, which gives the value
of the dependent variable for every combination of the independent variables.
A function of $n$ variables is defined on $2^{n}$ points, and its table has
$2^{n}$ rows, listed in counting order: 000, 001, 010, \ldots\ The function
$F$ of Table~\ref{tab:b2:F}, which is 1 on rows 3, 4, 5, 6 and 7, will be
followed through the whole chapter.

\begin{table}[htb]
\centering\small
\caption{The function $F$, with the minterm or the maxterm of each
row.}\label{tab:b2:F}
\begin{tabular}{@{}cccc|cll@{}}
\toprule
row & $A$ & $B$ & $C$ & $F$ & \textbf{minterm ($F=1$)} & \textbf{maxterm ($F=0$)}\\
\midrule
0 & 0&0&0 & 0 & & $M_0=A+B+C$\\
1 & 0&0&1 & 0 & & $M_1=A+B+\nott{C}$\\
2 & 0&1&0 & 0 & & $M_2=A+\nott{B}+C$\\
3 & 0&1&1 & 1 & $m_3=\nott{A}\,B\,C$ & \\
4 & 1&0&0 & 1 & $m_4=A\,\nott{B}\,\nott{C}$ & \\
5 & 1&0&1 & 1 & $m_5=A\,\nott{B}\,C$ & \\
6 & 1&1&0 & 1 & $m_6=A\,B\,\nott{C}$ & \\
7 & 1&1&1 & 1 & $m_7=A\,B\,C$ & \\
\bottomrule
\end{tabular}
\end{table}

\subsection{The half adder}

The \term{half adder} adds two bits $A$ and $B$ and produces a sum bit $S$ and
a carry $C$. In binary $1+1$ gives sum 0 and carry 1; from the table, $S$ is
1 when the inputs differ and $C$ is 1 when both are 1
(Figure~\ref{fig:b2:ha}):
\[ S = A\oplus B,\qquad C=A\cdot B. \]
The circuit is exactly the two instructions of the warm-up. It is ``half''
an adder because it has no carry input: it can only add the least significant
bits of two numbers.

\begin{figure}[htb]
\centering
\begin{minipage}{0.3\textwidth}\centering\small\ttfamily
\begin{tabular}{cc|cc}\textrm{$A$}&\textrm{$B$}&\textrm{$S$}&\textrm{$C$}\\\hline 0&0&0&0\\0&1&1&0\\1&0&1&0\\1&1&0&1\end{tabular}
\end{minipage}
\begin{minipage}{0.55\textwidth}\centering
\begin{circuitikz}
  \node[xor port] (x) at (3,1.2) {};
  \node[and port] (a) at (3,-0.6) {};
  \draw (x.in 1) -- ++(-1.6,0) node[left]{$A$} coordinate (A);
  \draw (x.in 2) -- ++(-1.6,0) node[left]{$B$} coordinate (B);
  \draw (x.in 1) ++(-1.0,0) node[circ]{} |- (a.in 1);
  \draw (x.in 2) ++(-0.5,0) node[circ]{} |- (a.in 2);
  \draw (x.out) -- ++(0.5,0) node[right]{$S$};
  \draw (a.out) -- ++(0.5,0) node[right]{$C$};
\end{circuitikz}
\end{minipage}
\caption{The half adder.}\label{fig:b2:ha}
\end{figure}

\subsection{The full adder}

The \term{full adder} adds two bits and a carry in, $C_{in}$, and produces a
sum bit and a carry out, $C_{out}$ (Figure~\ref{fig:b2:fa}). $S$ is 1 when an
odd number of inputs is 1; $C_{out}$ is 1 when at least two inputs are 1:
\[ S=A\oplus B\oplus C_{in},\qquad C_{out}=A\,B+C_{in}\,(A\oplus B). \]
The circuit is two half adders and an OR.

\begin{figure}[htb]
\centering
\begin{minipage}{0.32\textwidth}\centering\small\ttfamily
\begin{tabular}{ccc|cc}\textrm{$A$}&\textrm{$B$}&\textrm{$C_{in}$}&\textrm{$S$}&\textrm{$C_{out}$}\\\hline
0&0&0&0&0\\0&0&1&1&0\\0&1&0&1&0\\0&1&1&0&1\\1&0&0&1&0\\1&0&1&0&1\\1&1&0&0&1\\1&1&1&1&1\end{tabular}
\end{minipage}
\begin{minipage}{0.64\textwidth}\centering
\begin{circuitikz}[scale=0.9,transform shape]
  \node[xor port] (x1) at (2.4,2.0) {};
  \node[xor port] (x2) at (5.6,1.72) {};
  \node[and port] (a1) at (5.6,0.1) {};
  \node[and port] (a2) at (5.6,-1.5) {};
  \node[or port]  (o)  at (8.0,-0.7) {};
  \draw (x1.in 1) -- ++(-1.3,0) node[left]{$A$};
  \draw (x1.in 2) -- ++(-1.3,0) node[left]{$B$};
  \draw (x1.out) -- (x2.in 1);
  \draw (x2.in 2) -- ++(-0.6,0) coordinate (c) |- (-0.05,-2.6) node[left]{$C_{in}$};
  \draw (x1.out) ++(0.35,0) node[circ]{} |- (a1.in 1);
  \draw (c) ++(0,-1.0) node[circ]{} -- ++(0.0,0) |- (a1.in 2);
  \draw (x1.in 1) ++(-0.9,0) node[circ]{} |- (a2.in 1);
  \draw (x1.in 2) ++(-0.45,0) node[circ]{} |- (a2.in 2);
  \draw (a1.out) -| (o.in 1);
  \draw (a2.out) -| (o.in 2);
  \draw (x2.out) -- ++(2.6,0) node[right]{$S$};
  \draw (o.out) -- ++(0.3,0) node[right]{$C_{out}$};
\end{circuitikz}
\end{minipage}
\caption{The full adder: two half adders and an OR. Wires that cross without a
dot are not connected.}\label{fig:b2:fa}
\end{figure}

% ----------------------------------------------------------------------------
\section{Canonical forms}\label{sec:b2:canonical}

Canonical forms are a systematic way from the table to an expression.

\subsection{Elementary forms: functions of level 2}

\begin{definition}{Sum of products and product of sums}
An expression of \term{type P}, or \term{sum of products} (SOP), is a logical
sum of clauses, such as $y=a\,\nott{b}+\nott{a}\,d+b\,\nott{d}\,x$. It is
built as a layer of AND gates feeding one OR gate. An expression of
\term{type S}, or \term{product of sums} (POS), is a logical product of
factors, such as $y=(a+\nott{b})(\nott{a}+d)(b+\nott{d}+x)$. It is built as a
layer of OR gates feeding one AND gate.
\end{definition}

\subsection{Minterms and maxterms}

\begin{definition}{Minterm and maxterm}
A \term{minterm} of $n$ variables is a product that contains one literal of
every variable. A \term{maxterm} is a sum that contains one literal of every
variable.
\end{definition}

A minterm is 1 for exactly one row of the truth table, and a maxterm is 0 for
exactly one row. The minterm $m_i$ of row $i$ contains a variable
complemented where the row has a 0; the maxterm $M_i$ contains a variable
complemented where the row has a 1 (Table~\ref{tab:b2:F}). By De Morgan,
$M_i=\nott{m_i}$.

\subsection{From the truth table to the canonical forms}

\begin{theorem}{Canonical forms}
Every Boolean function is equal to
\begin{enumerate}[nosep]
\item the sum of the minterms of the rows where it is 1 (\term{canonical form P});
\item the product of the maxterms of the rows where it is 0 (\term{canonical form S}).
\end{enumerate}
Each of the two forms is unique, apart from the order of the terms.
\end{theorem}

\noindent\emph{Proof.} On a row where $f=1$, the minterm of that row is 1 and
so is the sum; on a row where $f=0$, all the minterms of the sum are 0. The
second statement is the dual. \hfill$\square$

\smallskip
For the function $F$:
\begin{align*}
F &= \nott{A}BC+A\nott{B}\,\nott{C}+A\nott{B}C+AB\nott{C}+ABC \;=\; \sum m(3,4,5,6,7),\\
F &= (A+B+C)\,(A+B+\nott{C})\,(A+\nott{B}+C) \;=\; \prod M(0,1,2).
\end{align*}

The theorem also proves that AND, OR and NOT are functionally complete. It is
synthesis by recipe: always possible, rarely economical.

\subsection{Algebraic minimisation of $F$}

The canonical form P of $F$ has 5 clauses and 15 literals. The algebra of
Chapter~\ref{ch:boole1} reduces it:
\begin{align*}
F &= \nott{A}BC+A\nott{B}\,\nott{C}+A\nott{B}C+AB\nott{C}+ABC\\
 &= \nott{A}BC+A\,(\nott{B}\,\nott{C}+\nott{B}C+B\nott{C}+BC) && \text{distributive: collect $A$}\\
 &= \nott{A}BC+A\,\bigl(\nott{B}\,(\nott{C}+C)+B\,(\nott{C}+C)\bigr) && \text{distributive, twice}\\
 &= \nott{A}BC+A\,(\nott{B}+B) && \text{complement: $\nott{C}+C=1$}\\
 &= \nott{A}BC+A && \text{complement: $\nott{B}+B=1$}\\
 &= A+BC && \text{simplification: $x+\nott{x}\,y=x+y$.}
\end{align*}
Fifteen literals have become three: from five AND gates with three inputs and
an OR with five, to one AND and one OR. It works, but each step needs an
idea. What is wanted is a method.

\subsection{The goal of minimisation}

The goal is to find, among all the possible two-level forms of type P, a
\emph{minimal} elementary form. Fewer and smaller gates mean lower cost, less
power and a faster circuit. Three methods are classical.

\begin{itemize}
\item \textbf{Karnaugh maps}: graphical, by groups of adjacent cells of maximum
  area; practical up to four or five variables \citep{karnaugh1953}.
\item \textbf{Quine's method}: algebraic, by systematic comparison of the terms.
\item \textbf{McCluskey's method}: tabular, suited to a program, and the basis
  of today's logic-synthesis tools. The Quine--McCluskey procedure is exact
  but exponential; industrial tools use heuristics.
\end{itemize}

% ----------------------------------------------------------------------------
\section{Karnaugh maps}

\subsection{The idea}

A Karnaugh map is the truth table of a function of $n$ variables redrawn on a
grid of $2^{n}$ cells, one cell per row. The cells are ordered so that two
neighbours differ in one variable only: hence the order 00, 01, 11, 10 of the
columns, which is a Gray code (Chapter~\ref{ch:combinational}), and hence the
fact that in a map $m_3$ comes before $m_2$. The left and right edges are
adjacent too: the map wraps round.

Graphical adjacency is thus logical adjacency: two adjacent 1s are two
minterms that differ in one variable, and by $x\,y+x\,\nott{y}=x$ they merge
into a clause with one literal fewer.

\begin{prgconcept}[The rule of the map]
Group the 1s in rectangles of 1, 2, 4 or 8 cells, as large as possible, until
every 1 is covered. A group of $2^{k}$ cells eliminates $k$ variables: the
clause of a group contains only the variables that do not change inside it.
\end{prgconcept}

\subsection{Two, three and four variables}

\begin{figure}[htb]
\centering
\begin{subfigure}[b]{0.25\textwidth}\centering
\resizebox{0.8\textwidth}{!}{\begin{karnaugh-map}[2][2][1][$B$][$A$]
  \minterms{1,2,3}\autoterms[0]
  \implicant{2}{3}\implicant{1}{3}
\end{karnaugh-map}}
\caption{$F=A+B$}\label{fig:b2:k2}
\end{subfigure}\hfill
\begin{subfigure}[b]{0.36\textwidth}\centering
\resizebox{\textwidth}{!}{\begin{karnaugh-map}[4][2][1][$BC$][$A$]
  \minterms{3,4,5,6,7}\autoterms[0]
  \implicant{4}{6}\implicant{3}{7}
\end{karnaugh-map}}
\caption{$F=A+BC$}\label{fig:b2:k3}
\end{subfigure}\hfill
\begin{subfigure}[b]{0.36\textwidth}\centering
\resizebox{\textwidth}{!}{\begin{karnaugh-map}[4][4][1][$CD$][$AB$]
  \minterms{0,2,5,7,8,10,13,15}\autoterms[0]
  \implicant{5}{15}\implicantcorner
\end{karnaugh-map}}
\caption{$F=BD+\nott{B}\,\nott{D}$}\label{fig:b2:k4}
\end{subfigure}
\caption{Karnaugh maps of two, three and four variables.}
\label{fig:b2:kmaps}
\end{figure}

\paragraph{Two variables (Figure~\ref{fig:b2:k2}).} The function is 1 for
$AB=01$, 10 and 11; its canonical form $\nott{A}B+A\nott{B}+AB$ has six
literals. One group is the row $A=1$: inside it $B$ changes, so $B$
disappears and the clause is $A$. The other is the column $B=1$: $A$ changes,
and the clause is $B$. Hence $F=A+B$: the function is simply OR. The cell
$AB$ belongs to both groups; using a 1 twice is the idempotence law,
$x+x=x$.

\paragraph{Three variables (Figure~\ref{fig:b2:k3}).} This is the function
$F$ of the chapter: a 1 in the cells of minterms 3, 4, 5, 6 and 7. A group of
four covers the whole row $A=1$: $B$ and $C$ both change, and the clause is
$A$. A group of two covers the column $BC=11$: $A$ changes, and the clause is
$BC$. Every 1 is covered and no group can be made larger: $F=A+BC$, the same
result as the algebra, found by eye in seconds.

\paragraph{Four variables (Figure~\ref{fig:b2:k4}).} The function
$F=\sum m(0,2,5,7,8,10,13,15)$ has eight minterms and thirty-two literals. A
group of four in the centre has $B=1$ and $D=1$ throughout: the clause is
$BD$. The four corners are also one group, because the map wraps round both
horizontally and vertically, as if it were drawn on a torus: there $B=0$ and
$D=0$, and the clause is $\nott{B}\,\nott{D}$. The result has four literals,
and can be recognised: $F$ is the XNOR of $B$ and $D$.

The practical rule is: always take the largest groups first, then cover what
is left.

\subsection{Don't-care conditions}

Sometimes a combination of the inputs can never occur, or the output does not
matter when it does. A digit in binary-coded decimal, for example, uses the
patterns 0000 to 1001, and the six patterns 1010 to 1111 never appear. Those
cells are marked in the map with a dash (or, in many texts, with a cross):
\term{don't care}. A don't-care cell
may be treated as a 1 when it makes a group larger, and as a 0 otherwise.
More freedom means larger groups and a cheaper circuit. The classic
exercise is the decoder for a seven-segment display, where each segment is a
function of four variables with six don't-cares.

% ----------------------------------------------------------------------------
\section{Back to the machine}

\subsection{Thirty-two full adders make an \texttt{add}}

The instruction \code{add t2, t0, t1} feeds the two registers into thirty-two
full adders, one per bit position (Figure~\ref{fig:b2:ripple}). The carry out
of each position is the carry in of the next: a \term{ripple-carry adder}, in
which the carry ripples leftwards as it does in a sum done by hand.

\begin{figure}[htb]
\centering
\begin{tikzpicture}[fa/.style={blkacc,minimum width=17mm,minimum height=12mm}]
  \foreach \i/\x/\l in {0/10.5/0, 1/7.9/1, 2/5.3/2, 3/1.3/31}{
    \node[fa] (f\i) at (\x,0) {full adder\\\footnotesize bit \l};
    \draw[<-] ([xshift=-3mm]f\i.north) -- ++(0,6mm) node[above,font=\ttfamily\small]{A\l};
    \draw[<-] ([xshift=3mm]f\i.north) -- ++(0,6mm) node[above,font=\ttfamily\small]{B\l};
    \draw[->] (f\i.south) -- ++(0,-6mm) node[below,font=\ttfamily\small]{S\l};
  }
  \draw[->,very thick,prgamber] (12.4,0) node[right,font=\ttfamily]{0} -- (f0.east);
  \draw[->,very thick,prgamber] (f0.west) -- (f1.east);
  \draw[->,very thick,prgamber] (f1.west) -- (f2.east);
  \draw[->,very thick,prgamber] (f2.west) -- ++(-0.8,0);
  \node[text=prgamber,font=\bfseries] at (3.3,0) {$\cdots$};
  \draw[->,very thick,prgamber] (3.0,0) -- (f3.east);
  \draw[->,very thick,prgamber] (f3.west) -- ++(-0.9,0) node[left,lbl,text=prgamber,align=right]{carry\\out};
  \node[lbl,text=prgamber] at (6.6,-0.9) {\emph{the carry ripples leftwards}};
\end{tikzpicture}
\caption{A ripple-carry adder of 32 bits.}\label{fig:b2:ripple}
\end{figure}

The carry out of bit 31 is the unsigned overflow of
Chapter~\ref{ch:integers}, which \riscv simply drops. Real processors use
faster schemes for the carry, such as carry look-ahead
(Chapter~\ref{ch:combinational}), but the function computed is the same.

\subsection{A full adder, one gate per instruction}

\asmfile[caption={The full adder as a program.},label=lst:b2:fa]{code/l09/c04.S}

Five instructions, five gates: Listing~\ref{lst:b2:fa} is the schematic of
Figure~\ref{fig:b2:fa}, read from left to right. Single-stepping it in \prgqt
and filling in the truth table is \emph{analysis}; deriving these gates from
the table with the method of this chapter was \emph{synthesis}. And with
whole registers as inputs, the same five lines work on all thirty-two bit
positions at once.

% ----------------------------------------------------------------------------
% Generated by tools/textbook/export_listings.js from decks/plan.js and code/l09.S: do not edit.
\section{The program of this chapter}\label{sec:lab:l09}

The listings of this chapter are parts of one program, \fname{code/l09.S} in the
repository of the course, printed in full in Listing~\ref{lst:lab:l09}. The
automatic tests of the course material build it and run it on the headless
runner of \prgqt{}, with scripted input, and check the frame;
Figure~\ref{fig:lab:l09} is the frame that the test wrote.

\begin{description}[style=nextline,leftmargin=1.5em]
\item[On the screen.] Eight rows, one per combination of A, B and Cin (000 at the top): three lamps for the inputs on the left, then S and Cout on the right. It is the truth table of the full adder.
\item[What to press.] Nothing: the program runs by itself.
\end{description}

\begin{figure}[htb]
\centering
\includegraphics[width=0.5\textwidth]{lab_l09.png}
\caption{\fname{l09.S} on \prgqt{}.}\label{fig:lab:l09}
\end{figure}

\asmfile[caption={The complete program of the chapter, \fname{code/l09.S}.},label=lst:lab:l09]{code/l09/program.S}

To run it on \prgqt{} the commands are those of Listing~\ref{lst:lab:l09:run};
Section~\ref{sec:prg:testing} gives the same steps for QEMU and for the board.

\begin{shell}[caption={Building \fname{l09.S} and running it on \prgqt{}.},label=lst:lab:l09:run]
# 1. start PRG32-QT on this computer (it listens on port 8080)

# 2. in the PRG32 folder, with ESP-IDF active: build the cartridge
python3 -m prg32 cartridge build ~/arc12/code/l09.S --portable \
    --entry-prefix l09 --name l09 --out l09.prg32

# 3. send it to PRG32-QT and start it
python3 -m prg32 esp32c6 upload-and-run l09.prg32 \
    --url http://localhost:8080 --slot cart0

# 4. look, press, compare with the previous slide
\end{shell}


% ----------------------------------------------------------------------------
\section{Check your understanding}

\begin{quiz}{How many distinct Boolean functions of three variables are there?}
\item 8
\item 9
\item 64
\item 256
\end{quiz}

\begin{quiz}{Which minterm of three variables $A$, $B$, $C$ corresponds to row 5 of the truth table?}
\item $\nott{A}\,B\,\nott{C}$
\item $A\,\nott{B}\,C$
\item $A+\nott{B}+C$
\item $A\,B\,\nott{C}$
\end{quiz}

\begin{quiz}{In a Karnaugh map, a rectangular group of four adjacent 1s eliminates:}
\item one variable
\item two variables
\item three variables
\item four variables
\end{quiz}

\begin{quiz}{In a full adder, the carry out is 1 when:}
\item exactly one input is 1
\item an odd number of inputs is 1
\item at least two inputs are 1
\item all three inputs are 1 only
\end{quiz}

\section{Exercises}

\begin{exercise}[paper first, then \prg]
\begin{enumerate}[nosep]
\item Write the canonical forms P and S of the outputs $S$ and $C_{out}$ of the
  full adder.
\item Minimise $C_{out}$ with a Karnaugh map of three variables. How many
  gates are saved?
\item Minimise $F=\sum m(0,1,2,5,8,9,10)$ with a map of four variables.
\item Build NOT, AND and OR using NAND only, and check each with a truth table.
\item Program the full adder on \prg and display $S$ and $C_{out}$ as lamps
  for the eight combinations of the inputs.
\end{enumerate}
\expected{On paper: the canonical form of $C_{out}$ has four minterms (3, 5,
6, 7) and the map gives three clauses of two literals,
$C_{out}=AB+AC_{in}+BC_{in}$, the \emph{majority} function. The third task
gives $F=\nott{B}\,\nott{D}+\nott{B}\,\nott{C}+\nott{A}\,\nott{C}\,D$. On
\prg: eight pairs of lamps that reproduce the truth table of the full adder.}
\end{exercise}

\section{Glossary}

\begin{glossarytable}
literal & letterale\\
clause, product term & clausola, termine prodotto\\
factor, sum term & fattore, termine somma\\
minterm / maxterm & mintermine / maxtermine\\
half adder / full adder & semisommatore / sommatore completo\\
carry in / carry out & riporto entrante / uscente\\
canonical form & forma canonica\\
sum of products (SOP) & somma di prodotti (tipo P)\\
product of sums (POS) & prodotto di somme (tipo S)\\
Karnaugh map & mappa di Karnaugh\\
don't care & condizione di indifferenza\\
ripple-carry adder & sommatore a propagazione del riporto\\
\end{glossarytable}

\section{Conclusions and reading}

A Boolean function is a truth table, and a truth table always gives a
canonical expression. Minimisation is the process that optimises a logic
network in order to simplify its implementation. Karnaugh maps turn the
algebra into a picture: group the 1s, read the clauses. The half adder and
the full adder are the first real circuits of this book, and thirty-two full
adders in a row are what executes the \ins{add} typed in
Chapter~\ref{ch:firstcontact}.

\paragraph{Reading.} \citet{patterson2020}, appendix~A; \citet{bucci2017},
the chapters on combinational networks; Charles Petzold, \emph{Code}, the
chapters on logic gates and on binary adding machines.
